\documentclass{ctexart}
\usepackage{geometry}
\usepackage{fancyhdr}
\usepackage{graphicx}
\usepackage{booktabs}
\usepackage{amsmath}
\usepackage{tikz}
\usepackage{array}
\xeCJKsetup{CJKmath=true} 
\usepackage{zhnumber} % change section number to chinese
\renewcommand\thesection{\zhnum{section}}
\renewcommand \thesubsection {\arabic{subsection}}
\CTEXsetup[format={\Large\bfseries}]{section}

\geometry{
    a4paper,
    left=3.18cm,
    right=3.18cm,
    top=3.04cm,
    bottom=3.04cm
}

\pagestyle{fancy}
\fancyhf{}
\renewcommand{\headrulewidth}{0.7pt} % 设置页眉横线粗细
\fancyhead[L]{\kaishu\large 大学物理实验报告} % 在左侧设置页眉文字
\fancyhead[R]{\kaishu\large 哈尔滨工业大学(深圳) } % 在右侧设置页眉文字
\fancyfoot[R]{\thepage} % 将页数放在右下角


\setlength\headwidth{\textwidth}

\begin{document}

\noindent
\begin{center}
\textbf{
\begin{tabular}{p{2.4cm}p{2.4cm}p{4cm}p{4cm}}
    班级 \hrulefill & 学号 \hrulefill & 姓名 \hrulefill & 教师签字 \hrulefill \\
\end{tabular}
\begin{tabular}{p{6cm}p{3.6cm}p{3.6cm}}
    实验日期 \hrulefill & 预习成绩 \hrulefill & 总成绩 \hrulefill
\end{tabular}
{\noindent}	 \rule[-10pt]{\textwidth}{0.7pt}
}\end{center}

\begin{center}
    \Large \textbf{实验内容 \underline{液体表面张力系数测量}}
\end{center}

\section{实验预习}
\subsection{什么是表面张力? 液体表面张力系数与哪些因素有关?}
\subsection{拉脱法测量液体表面张力的实验原理是什么？}

\newpage
\section{实验现象及原始数据记录}
\subsection{吊环的内、外直径（单位：mm）}

\begin{table}[h]
    \renewcommand\arraystretch{1.6}
    \centering
    \caption{吊环的测量}
    \label{tab:m}
    \begin{tabular}{|m{1.8cm}<{\centering}|m{1.5cm}<{\centering}|m{1.5cm}<{\centering}|m{1.5cm}<{\centering}|m{1.5cm}<{\centering}|m{1.5cm}<{\centering}|m{1.5cm}<{\centering}|}
        \hline
        测量次数& 1 & 2 & 3 & 4 & 5 & 平均值 \\
        \hline
        内径$D_i$ & & & & & & \\
        \hline
        外径$D_o$ & & & & & & \\
        \hline
    \end{tabular}
\end{table}

~

\subsection{利用逐差法求仪器的转换系数$K$：}
~\\
先记录砝码盘等作为初始读数$V_0 = \rule{1.3cm}{0.4pt} mV $ ，然后每次增加一个砝码500mg，（该标准砝码符合国家标准，相对误差为0.05\%）
~\\

\begin{table}[h]
    \renewcommand\arraystretch{1.3}
    \centering
    \caption{吊环的测量}
    \label{tab:k}
    \begin{tabular}{|m{2cm}<{\centering}|m{2cm}<{\centering}|m{2cm}<{\centering}|m{3.5cm}<{\centering}|}
        \hline
        砝码质量 $10^{-6} Kg $ & 增重读数 $V_1' (mV)$ & 减重读数 $V_1'' (mV)$ & $V_i = \frac{V_i{'}+V_i{''}}{2} (mV)$ \\
        \hline
        0.00 & & & \\
        \hline
        500.00 & & & \\
        \hline
        1000.00 & & & \\
        \hline
        1500.00 & & & \\
        \hline
        2000.00 & & & \\
        \hline
        2500.00 & & & \\
        \hline
        3000.00 & & & \\
        \hline
        3500.00 & & & \\
        \hline
    \end{tabular}
\end{table}

~

利用逐差法求出每 500 mg 对应的电子秤的读数 $\Delta V$ ，则 $\overline{K} = \frac{mg}{\Delta V} = \rule{1.5cm}{0.4pt} N/mV $ 
\newpage
\subsection{用拉脱法求拉力对应的电子秤读数：}

\begin{table}[h]
    \renewcommand\arraystretch{1.6}
    \centering
    \caption{室温下表面张力系数测量表}
    \label{tab:T1}
    \begin{tabular}{|m{1.8cm}<{\centering}|m{3cm}<{\centering}|m{2.5cm}<{\centering}|m{4cm}<{\centering}|}
        \multicolumn{4}{c}{水温（室温）$\rule{1.5cm}{0.4pt} ^\circ C$,  电子秤初始读数$V_0 = \rule{1.5cm}{0.4pt} mV$} \\
        \hline
        测量次数& 拉脱时最大读数 $V_1 (mV)$ & 吊环读数 $V_2 (mV)$ & 表面张力对应读数(mV) $V_i^\ast = V_1-V_2$\\
        \hline
        1 & & & \\
        \hline
        2 & & & \\
        \hline
        3 & & & \\
        \hline
        4 & & & \\
        \hline
        5 & & & \\
        \hline
        平均值 & & & $\overline{V} = \rule{2.5cm}{0.4pt} $\\
        \hline
    \end{tabular}
\end{table}

\begin{table}[!h]
    \renewcommand\arraystretch{1.6}
    \centering
    \caption{不同温度下表面张力系数测量表}
    \label{tab:T2}
    \begin{tabular}{|m{1.8cm}<{\centering}|m{3cm}<{\centering}|m{2.5cm}<{\centering}|m{4cm}<{\centering}|}
        \multicolumn{4}{c}{水温（室温）$\rule{1.5cm}{0.4pt} ^\circ C$,  电子秤初始读数$V_0 = \rule{1.5cm}{0.4pt} mV$} \\
        \hline
        测量次数& 拉脱时最大读数 $V_1 (mV)$ & 吊环读数 $V_2 (mV)$ & 表面张力对应读数(mV) $V_i^\ast = V_1-V_2$\\
        \hline
        1 & & & \\
        \hline
        2 & & & \\
        \hline
        3 & & & \\
        \hline
        4 & & & \\
        \hline
        5 & & & \\
        \hline
        平均值 & & & $\overline{V} = \rule{2.5cm}{0.4pt} $\\
        \hline
    \end{tabular}
\end{table}

\begin{tikzpicture}[remember picture,overlay]
    \node[anchor=south east,inner sep=100pt] at (current page.south east) {
        \renewcommand{\arraystretch}{1.5} % 表格行高倍数
        \setlength{\tabcolsep}{18pt}    
    \begin{tabular}{|c|c|}
        \hline
        \LARGE  教师 & \LARGE  姓名 \\
        \hline
        \LARGE \kaishu 签字 &  \\
        \hline
        \end{tabular}
    };
\end{tikzpicture}

\newpage

\section{数据处理}
\subsection{测量室温下水的表面张力系数，并计算不确定度}

\begin{align*}
    \overline{L} &= \pi \left(\overline{D_{\text{内}}} + \overline{D_{\text{外}}}\right) \\
    \overline{\alpha} &= \frac{\overline{K} \cdot \overline{V}}{\overline{L}} \\
    \left(\frac{\Delta \alpha}{\alpha}\right)^2 &= \left(\frac{\Delta K}{K}\right)^2 + \left(\frac{\Delta V}{V}\right)^2 + \left(\frac{\Delta L}{L}\right)^2 \\
    \alpha &= \alpha \pm \Delta \alpha 
\end{align*}

将室温下的$D_{\text{内}}$和$D_{\text{外}}$代入第一个公式，计算得：

$$ \overline{L} = \pi \left(\overline{D_{\text{内}}} + \overline{D_{\text{外}}}\right) = \pi \left(33.08 + 34.81\right) = 213.17 mm $$

然后使用逐差法计算$\Delta V$，进而计算$\overline{K}$和$\overline{\alpha}$：

$$ \Delta V = \frac{\sum_{i = 5}^{8} V_i - \sum_{i = 1}^{4} V_i}{16} = 0.491 mV $$

$$ \overline{K} = \frac{mg}{\Delta V} = \frac{0.5 \times 10^{-3} \times 9.8}{\Delta V} = 9.98 \times 10^{-3} N/mV $$

$$ \overline{\alpha} = \frac{\overline{K} \cdot \overline{V}}{\overline{L}} = \frac{9.98 \times 10^{-3} \times 1.661}{213.17 \times 10^{-3}} = 7.78 \times 10^{-2} N/m $$

下面计算不确定度。已知实验中所用的500mg的砝码精度为0.25mg，电子秤最小读数为0.001mV，所以有：

\begin{align*}
\left(\frac{\Delta K}{K}\right)^2 &= \left(\frac{\Delta V}{V}\right)^2 + \left(\frac{\Delta m}{m}\right)^2 \\
&= \left(\frac{\frac{0.001}{\sqrt{3}}mV}{\frac{1}{16}\times7.856mV}\right)^2 + \left(\frac{\frac{0.25}{\sqrt{3}}mg}{500mg}\right)^2 \\
&\approx 1.466 \times 10^{-6}
\end{align*}

而对于主不确定度中$\left(\frac{\Delta V}{V}\right)^2$的值，其计算方式为：

$$ \left(\frac{\Delta V}{V}\right)^2 = \left(\frac{\frac{0.001}{\sqrt{3}}mV}{1.661mV}\right)^2 = 1.208 \times 10^{-7} $$

而20分度的游标卡尺的仪器误差为0.05mm，所以有：

$$ \left(\frac{\Delta L}{L}\right)^2 = \left(\frac{\frac{0.05}{\sqrt{3}}mm}{\overline{D_\text{内}}+\overline{D_\text{外}}}\right)^2 \approx 1.808 \times 10^{-7} $$

将以上结果代入第四个公式，计算得：

$$ \left(\frac{\Delta \alpha}{\alpha}\right)^2 = \left(\frac{\Delta K}{K}\right)^2 + \left(\frac{\Delta V}{V}\right)^2 + \left(\frac{\Delta L}{L}\right)^2 = 1.768 \times 10^{-6} $$

则不确定度为：

$$ \Delta\alpha = \alpha\times\sqrt{1.768 \times 10^{-6}} = 7.78 \times 10^{-2} \times 1.33 \times 10^{-3} \approx 1.03 \times 10^{-4} \, \text{N/m} $$

故最终表面张力系数为：

$$ \alpha = \overline{\alpha} \pm \Delta\alpha = (7.78 \pm 0.10) \times 10^{-2} \, \text{N/m} $$

置信概率$P = 68.3\%$

\subsection{从附录中查出室温下水的表面张力系数$\alpha$的理论值，把实验结果与此值比较求相对误差,并进行分析。}

从附录中查得在$25^\circ C$下，纯净水表面张力系数的公认值为$7.197\times10^{-2}N/m$，故相对误差为：

$$ \sigma = \frac{\vert 7.78-7.197\vert }{7.197} = 8.1\% $$

此误差不大，可以接受。误差可能来自吊环底面与水面不平行，同时也可能有来自拉断液膜时读数不及时造成的影响。

\subsection{测量不同温度下表面张力系数，并与室温下水的表面张力系数理论值作分析比较。}

计算过程与上述计算过程相同，省略计算过程。计算$40.5^\circ C$下表面张力系数得：

$$ \alpha = (6.95 \pm 0.09) \times 10^{-2} N/m $$
$$ P = 68.3 \% $$

经查表得到在$40^\circ C$下水的表面张力系数为$6.956\times 10^{-2} N/m$，所以相对误差为：

$$ \sigma = \frac{\vert 6.956-6.95 \vert}{6.956} = 0.09\% $$

经计算得到$40.5^\circ C$下其表面张力系数比$26.7^\circ C$要小一些，这是因为温度升高，分子热运动加剧，液体分子之间的吸引力会减小，于是表面张力系数减小。

\section{实验结论及现象分析}
（讨论液体表面张力系数测量中误差的来源，如何提高测量精度？）

实验结果表明，测量值与理论值稍有偏差，但是在实验误差允许范围内，可以接受。

液体表面张力系数测量的误差来自多个方面。首先，测量仪器转换系数时，若未等到砝码盘稳定后再读数，则由于加速度的影响会出现误差；其次，测量金属环直径时存在游标卡尺读数的误差；再次，吊环底面与水面不平行也会导致误差；最后，调整水位过快可能会导致拉断液面时测力计还未来得及探测到峰值，产生读数误差。

为了提高测量精度，首先需要保持待测液体的纯净和金属环、砝码的清洁。其次，在操作时，要保持金属环与水面平行，同时避免实验环境出现较大的振动。最后，使用精度更高的测力计和游标卡尺可以提高读数精度。

\section{讨论题}
\subsection{在推导液体表面张力系数测量公式中作了哪些近似？式中各量的物理意义是什么？}

所做的近似有：1.湿润角近似为零。2.认为拉起的液膜质量可忽略。

$$\alpha = \frac{F-m_0g}{\pi (D_{\text{内}}+D_{\text{外}})}$$

上式中，$F$为拉起吊环的拉力，$m_0g$为吊环的重力。$D_{\text{内}}，D_{\text{外}}$为吊环的内、外径。

\subsection{若考虑拉起液膜的重量，实验结果应如何修正？}
考虑拉起液膜的重量，计算公式应修正为：

$$\alpha = \frac{F-(m+m_0)g}{\pi (D_{\text{内}}+D_{\text{外}})}$$

若能求出液膜的质量，则可以代入新的公式重新计算来对实验结果进行修正。 
\end{document}